% analysis/part_intro-DE.tex
\chapter*{Analyse}
\phantomsection
\addcontentsline{toc}{chapter}{Analyse}

Dieser Teil legt dar, wie AIM als kohärentes Verarbeitungssystem und nicht als
unverbundene Sammlung von Modellen arbeitet.

\section*{Was wissen wir bereits?}
Die vorangehenden Teile haben die begrifflichen, geometrischen, operativen,
realisierungsbewussten und modellierungstechnischen Grundlagen von AIM
geschaffen.

Eisenbahn-Alignments können nun durch dünn besetzte semantische Modelle
dargestellt, durch Laufzeitdienste rekonstruiert, durch Zustandsoperatoren
interpretiert und durch Realisierungsbegriffe mit der physischen Realität in
Beziehung gesetzt werden.

Die einzelnen Bausteine des Rahmens sind damit vorhanden. Eine Sammlung von
Komponenten bildet jedoch noch kein vollständiges System.

Das System besitzt außerdem eine Entwicklungslinie, die erklärt, weshalb diese
Komponenten existieren. Berlinish begann als ursprüngliche Forschungsgrammatik
für Übergangsfunktionen. axtranNew sollte diese Grammatik in lösbare
Alignment-Berechnungen überführen. transitionDB und TransEd behandeln die
Bewahrung und Untersuchung des wiederverwendbaren Funktionswissens. SPOT und
alignmentOS behandeln konkrete Objektidentität, Zustand, Bearbeitung und
operative Nutzung. Der Knowledge Kernel stellt die stabile Terminologie und
die Verantwortungsgrenzen bereit, welche die ursprüngliche Berechnungsarbeit
allein nicht liefern konnte. Research stellt Aussagen infrage und erweitert
sie; die Thesis erklärt, verbindet und überträgt sie.

Diese Verantwortungen bilden eine Ausführungs- und Wissenskette. Sie dürfen
nicht als Behauptung gelesen werden, jede benannte Komponente sei ein
kanonischer Kernel-Begriff oder eine vollständige Implementierung.

\section*{Was fehlt noch?}
Praktische Eisenbahninformationsmodellierung erfordert mehr als isolierte
Modelle und Operatoren. Daten müssen Verarbeitungsketten durchlaufen, Zustände
transformiert werden, Dienste zusammenwirken, Fehler erkannt und Ergebnisse
aus Zwischendarstellungen erzeugt werden.

Es fehlt daher die Verarbeitungsarchitektur, die einzelne Begriffe zu einem
kohärenten operativen Rahmen verbindet. AIM ist nicht nur als Sammlung von
Modellen, sondern als Pipeline zusammenwirkender Transformationen zu verstehen.

\section*{Was folgt?}
Dieser Teil führt die Systemsicht von AIM ein. Der Schwerpunkt verschiebt sich
von einzelnen Begriffen zu ihrem Zusammenwirken.

Die folgenden Kapitel untersuchen:
\begin{itemize}
\item Transformationspipelines,
\item Laufzeitdienste,
\item Operatorketten,
\item Systeminteraktionen,
\item Fehlermodi und
\item architektonische Entwurfsprinzipien.
\end{itemize}

Ziel ist zu erklären, wie Eisenbahninformation in AIM von Erfassung über
Bewertung und Realisierung bis zur Optimierung fließt. Diese Sicht zeigt AIM
als operatorbasiertes Eisenbahninformationssystem und nicht als bloße Sammlung
geometrischer Verfahren.

\section*{Bezug zur AIM-Roadmap}
In der AIM-Roadmap entspricht dieser Teil dem Übergang von Alignment-Modellen
zu Verarbeitungspipelines und Systemdiensten.

Die folgende Skizze beantwortet die Leitfrage: Wie schreitet das Argument von
der Modellstruktur zum operativen Systemverhalten fort?

\begin{center}
\begin{tikzpicture}[node distance=1.4cm,box/.style={draw,rounded corners,minimum width=5.0cm,minimum height=0.9cm,align=center}]
\node[box] (model) {Alignment-Modell};
\node[box, below=of model] (pipeline) {Pipelines und Dienste};
\node[box, below=of pipeline] (system) {Systemsicht};
\draw[->,thick] (model)--(pipeline); \draw[->,thick] (pipeline)--(system);
\end{tikzpicture}
\end{center}

Die Abfolge
\[
\text{Alignment-Modell}\rightarrow\text{Pipelines und Dienste}
\rightarrow\text{Systemsicht}
\]
markiert den Übergang von statischen Eisenbahndarstellungen zu operativen
Verarbeitungsarchitekturen. Diese Systemsicht bereitet Optimierung,
Engineering-Abläufe und praktische Bereitstellung vor.

\begin{principle}[Pipeline-Prinzip]
\label{princ:pipeline-principle}
In AIM wird Eisenbahninformation durch ausdrückliche Operatorpipelines
verarbeitet, in denen Zustände, Modelle, Realisierungen und Bewertungen durch
wohldefinierte Transformationen verbunden sind.
\end{principle}

\begin{summary}
Modellierung erklärt, wie Eisenbahninformation dargestellt wird.

Analyse erklärt, wie Eisenbahninformation verarbeitet wird.

Dieser Teil begründet damit die Systemarchitektur, die Modelle, Operatoren,
Dienste und Transformationen zu einem kohärenten AIM-Rahmen verbindet.
\end{summary}

Diese Systemsicht bereitet den Teil Optimierung vor, in dem
Pipeline-Ergebnisse für strukturierte Verbesserung und
Entscheidungsunterstützung verwendet werden.
